\documentclass[11pt,a4paper]{article}

\usepackage[numbers,sort&compress]{natbib}
\usepackage{amsmath,amssymb}
\usepackage{graphicx}
\usepackage[export]{adjustbox}
\usepackage{capt-of}
\usepackage{booktabs}
\usepackage{array}
\usepackage{tabularx}
\usepackage{multirow}
\usepackage[margin=2.2cm]{geometry}
\usepackage[colorlinks=true,linkcolor=black,citecolor=black,urlcolor=blue]{hyperref}

\graphicspath{{figures/}{./figures/}{./}{paper/figures/}}
\DeclareGraphicsExtensions{.pdf,.png,.jpg,.jpeg}

\newcolumntype{Y}{>{\raggedright\arraybackslash}X}
\newcolumntype{Z}[1]{>{\centering\arraybackslash}p{#1}}

\renewcommand{\thesection}{S\arabic{section}}
\renewcommand{\thesubsection}{S\arabic{section}.\arabic{subsection}}
\renewcommand{\thefigure}{S\arabic{figure}}
\renewcommand{\thetable}{S\arabic{table}}
\renewcommand{\theequation}{S\arabic{equation}}

\newcommand{\paperfig}[5][]{%
  \begin{center}
  \begin{minipage}{\linewidth}
  \centering
  \includegraphics[
    width=#2,
    max height=0.72\textheight,
    keepaspectratio,
    #1
  ]{#3}%
  \vspace{-0.35em}
  \captionof{figure}{#4}
  \label{#5}
  \end{minipage}
  \vspace{-0.45em}
  \end{center}
}

\newenvironment{suppTableBlock}[2]{%
  \begin{center}
  \begin{minipage}{\linewidth}
  \captionof{table}{#1}
  \label{#2}
  \centering
}{%
  \end{minipage}
  \vspace{0.35em}
  \end{center}
}

\title{Adaptation-Aware Model Selection for Few-Shot Digital Twin
Instantiation}
\author{}
\date{}

\begin{document}
\maketitle

This supplementary material provides additional implementation, analytical,
reproducibility, and diagnostic details for the main manuscript, including
counting conventions, candidate-bank and calibration records, analytical
derivations, baseline configurations, hardware profiling, real-workload
evaluations, and synthetic-benchmark generation settings.

The source, architecture-preselection, margin-calibration, diagnostic, and
held-out target pools use disjoint center identities.

\section{Detailed Operation-Count and Parameter-Count Rules}
\label{supp:counting}

The estimated operation count is computed for one forward pass with batch
size one. Additions and multiplications in the main linear, convolutional,
recurrent-matrix, residual-projection, and output-layer operations are
counted. Bias additions, nonlinear activations, dropout operations, and GRU
elementwise operations are excluded. The same deterministic rule is applied
to every candidate.

For an instantiated architecture $a_q$, the parameter count is the exact
number of elements in all trainable parameter tensors after the target input
dimension and prediction horizon are fixed. Both counts depend on the target
case only through its instantiated input and output dimensions. Because
parameter adaptation does not change the architecture, both quantities remain
unchanged.

The two indicators are constrained by hard deployment limits. They are not
added to the prediction loss and are not changed after target adaptation.
They are deterministic model-level indicators rather than direct guarantees
of hardware-specific latency or memory use.
Their relation to measured inference latency and serialized model size is
evaluated in Supplementary Section~S5.

\section{Candidate-Bank Records and Relative-Margin Calibration}
\label{supp:bank-calibration}

\subsection{Candidate-bank construction and source-trained records}

The initial architecture space contains 66 indexed configurations: 18 MLPs,
36 TCNs, and 12 GRUs. Configuration A57, a one-layer GRU with 32 hidden
units, is fixed as the reference architecture before preselection.

An alternative configuration is retained when it records at least two
check-oracle wins over A57, or at least two validation-selected positive check
wins over A57. A check-oracle win is assigned to the feasible configuration
with the lowest check loss when it improves on A57; a validation-selected
positive check win is assigned when validation selects an alternative
configuration whose check loss improves on A57.
These cases belong to the architecture-preselection pool and are disjoint from
the margin-calibration and held-out target pools.
The recorded S0 retained-indexed-configuration list contains three MLP and
three GRU entries, including A57. No TCN configuration satisfies this rule.
The PT-A57 reference candidate is constructed separately in the downstream
candidate bank.

Two retained one-layer GRU-32 configurations have different indexed
settings, but recurrent dropout is inactive for a one-layer GRU. They
therefore represent the same executable architecture. The six retained
configurations correspond to five executable architectures.

The final bank contains the pooled-source PT+FT initialization for
the reference, one source-trained initialization for each of the five other
retained indexed configurations, and a second A57 initialization used as a
paired-initialization control. The six retained configurations therefore
produce seven candidates. The bank and reference candidate are fixed
before relative-margin calibration and held-out evaluation.

Source preparation is not identical across the candidate records. Strong-bank
alternatives use 50 epochs of Adam/MSE training with learning rate $10^{-3}$
and batch size 64. The pooled-source PT-A57 reference is loaded from a
separately frozen asset whose manifest records a nominal 50-epoch source
recipe. The C1-A57 control has a distinct source-training lineage; its
released source builder specifies 10 Adam/MSE epochs followed by 180
Adam/Huber refinement updates. The historical per-step C1 execution trace is
not fully recoverable. Thus, the two A57 initializations differ in source
preparation as well as initialization, and the controls do not isolate a
random-seed-only or architecture-only effect.

\subsection{Relative-margin calibration records}

Section~4.4 of the main manuscript defines the calibration rule. This section
records its complete experimental instantiation and reproducibility details.

Let
$\mathcal{T}=\{0.05,0.075,0.10,0.125,0.15,0.20\}$
denote the preset margin grid. For a relative margin $\tau$, let
$\widehat R_{\mathrm{harm}}(\tau)$ denote the fraction of all calibration
cases in which an alternative is selected and its check-set loss exceeds that
of the adapted reference. Let $\widehat G_{\mathrm{ref}}(\tau)$ and
$\widehat G_{\mathrm{pair}}(\tau)$ denote the mean relative check-loss gains
with respect to the adapted reference and the paired-initialization control,
respectively, where each gain is the relative reduction in check-set loss of
the selected model with respect to the corresponding control. Their lower
95\% confidence bounds are
$\underline G_{\mathrm{ref}}(\tau)$ and
$\underline G_{\mathrm{pair}}(\tau)$.

The eligible margins form
\begin{equation}
\begin{aligned}
\mathcal{T}_{\mathrm{elig}}
=
\bigl\{
\tau\in\mathcal{T}
\;\bigm|\;&
\widehat R_{\mathrm{harm}}(\tau)\leq\rho,
\\
&
\widehat G_{\mathrm{ref}}(\tau)\geq\gamma_{\mathrm{ref}},
\quad
\underline G_{\mathrm{ref}}(\tau)>0,
\\
&
\widehat G_{\mathrm{pair}}(\tau)\geq\gamma_{\mathrm{pair}},
\quad
\underline G_{\mathrm{pair}}(\tau)>0
\bigr\}.
\end{aligned}
\label{eq:supp-margin-calibration}
\end{equation}

The calibration uses
$\rho=0.05$ and
$\gamma_{\mathrm{ref}}=\gamma_{\mathrm{pair}}=0.03$.
When $\mathcal{T}_{\mathrm{elig}}$ is nonempty, the calibrated margin is
\begin{equation}
\tau^{*}=\min \mathcal{T}_{\mathrm{elig}}.
\end{equation}
The smallest eligible value in the main margin-calibration setting is
$\tau^{*}=0.10$.
The paired-initialization control is a two-candidate selector, not the
unconditionally deployed second A57 model: both A57 source initializations
are adapted on the same support set, then the reference-relative validation
rule at the specified $\tau$ selects one for check evaluation.
The $\widehat G_{\mathrm{pair}}$ denominator is the check loss of this
paired-selector output in each corresponding calibration case.

Architecture preselection and candidate-bank construction are completed before
margin calibration, and the resulting candidate bank and reference candidate
remain fixed while the margin grid is evaluated.

If no margin satisfies the calibration criteria in a new domain, the adapted
reference is retained. A margin fixed in another domain may still be
evaluated without recalibration to assess transfer. The six-margin calibration results are given in
Table~\ref{tab:supp-synthetic-calibration}. The figures are from the frozen
20-center, 80-case check-set calibration records; each center contributes
four $H/K$ cases. Cluster-resampled 95\% confidence intervals use 4,000
repetitions. This is an empirical eligibility check, not a distribution-free
risk guarantee.

\begin{suppTableBlock}{Five-criterion margin calibration on the main synthetic development pool (20 centers, 80 cases).}{tab:supp-synthetic-calibration}
\footnotesize
\setlength{\tabcolsep}{4.2pt}
\renewcommand{\arraystretch}{1.12}
\begin{tabular}{lrrrrrl}
\toprule
$\tau$ & Harm (\%) & Ref. gain (\%) & Ref. LCB (\%) & Pair gain (\%) & Pair LCB (\%) & Eligible\\
\midrule
0.050 & 6.25 & 13.93 & 7.75 & 13.33 & 7.73 & No\\
0.075 & 6.25 & 13.79 & 7.77 & 13.18 & 7.39 & No\\
0.100 & 5.00 & 13.85 & 8.06 & 13.32 & 7.68 & Yes\\
0.125 & 5.00 & 13.13 & 7.01 & 12.60 & 6.85 & Yes\\
0.150 & 5.00 & 12.70 & 6.48 & 12.17 & 6.30 & Yes\\
0.200 & 2.50 & 10.98 & 5.56 & 10.45 & 5.70 & Yes\\
\bottomrule
\end{tabular}
\vspace{0.25em}
\par\footnotesize\raggedright
Harm is the fraction of all calibration cases with a harmful alternative
selected on the independent check set. Ref. and pair gains are mean
case-wise relative check-MSE reductions against the adapted PT reference
and the two-A57 paired-initialization selector, respectively; LCB denotes
the lower bound of the corresponding 95\% center-cluster bootstrap CI.
Eligibility requires Harm $\le5\%$, both mean gains $\ge3\%$, and both
lower bounds $>0$. The two smaller margins fail only the Harm criterion;
$\tau=0.10$ is the smallest eligible value. The exact frozen statistics
are provided in the repository's C3.2 selector calibration manifest.
\end{suppTableBlock}

\section{Additional Analytical Details}
\label{supp:analysis}

\subsection{Ideal Candidate Decision within the Adapted Feasible Set}

Let the expected future loss of adapted feasible candidate $q$ on target
center $c$ be
\begin{equation}
R_c^{\mathrm{cand}}(q)
=
R_c\left(a_q,\theta_{c,q}^{(T)}\right).
\label{eq:supp-candidate-risk}
\end{equation}
If this quantity were available, the ideal candidate-level decision would
satisfy
\begin{equation}
q_c^{\dagger}
\in
\underset{q\in\mathcal{Q}_{c}^{\mathrm{fea}}}
{\arg\min}\;
R_c^{\mathrm{cand}}(q).
\label{eq:supp-ideal-target}
\end{equation}

This candidate-level quantity is unavailable during instantiation. MSA-DTI
therefore uses post-adaptation validation loss $J_c(q)$ together with the
relative-margin rule in the main manuscript.

\subsection{Conditional expected-loss result}

Assume that validation loss differs from expected future loss by at most
$\epsilon_c\geq0$ for every feasible candidate:
\begin{equation}
\left|
R_c^{\mathrm{cand}}(q)-J_c(q)
\right|
\leq
\epsilon_c,
\qquad
\forall q\in\mathcal{Q}_{c}^{\mathrm{fea}}.
\label{eq:supp-validation-error}
\end{equation}

When MSA-DTI selects an alternative, the relative-margin rule gives
\begin{equation}
R_c^{\mathrm{cand}}(q_c^{*})
-
R_c^{\mathrm{cand}}(q^{\mathrm{ref}})
\leq
-\tau J_c(q^{\mathrm{ref}})
+
2\epsilon_c.
\label{eq:supp-expected-loss}
\end{equation}

Hence, the selected alternative has lower expected future loss than the
reference when
\begin{equation}
\tau J_c(q^{\mathrm{ref}})>2\epsilon_c.
\label{eq:supp-sufficient-condition}
\end{equation}
Equation~\eqref{eq:supp-sufficient-condition} is a sufficient condition under
assumption~\eqref{eq:supp-validation-error}. The present experiments do not
estimate an empirical bound on $\epsilon_c$.

\section{Additional Experimental and Diagnostic Results}
\label{supp:diagnostics}



\subsection{Target-side baseline settings}

For the synthetic main study, the window counts in support, validation,
check, and test order are $(10,1351,628,627)$, $(20,1346,625,625)$,
$(10,1346,624,624)$, and $(20,1342,621,621)$ for $(H,K)=(1,10)$,
$(1,20)$, $(4,10)$, and $(4,20)$, respectively.

Table~\ref{tab:supp-baselines} reports the target-side settings used in the
main comparison.

\begin{suppTableBlock}{Target-side settings of MSA-DTI and the baselines in the main comparison.}{tab:supp-baselines}
\footnotesize
\renewcommand{\arraystretch}{1.08}
\setlength{\tabcolsep}{2.0pt}
\begin{tabularx}{\linewidth}{
>{\raggedright\arraybackslash}p{0.18\linewidth}
>{\raggedright\arraybackslash}p{0.16\linewidth}
Y
>{\raggedright\arraybackslash}p{0.13\linewidth}
Z{0.08\linewidth}
>{\raggedright\arraybackslash}p{0.14\linewidth}}
\toprule
Method & Structure & Source initialization & Target update &
Adapted models & Deployment-limit check \\
\midrule
PT+FT & Fixed reference & Pooled-source model & SGD/MSE--50 & 1 &
Before output \\
MeDeT-based adaptation & Fixed reference & Meta-learned initialization &
SGD/MSE--50 & 1 & Before output \\
Random initialization & Fixed reference & Random & SGD/MSE--50 & 1 &
Before output \\
Few-shot NAS & Top-12 shortlist & Method-specific models & Adam/Huber--50 &
$\leq12$ & Before evaluation \\
Zero-shot NAS & Top-ranked model & Method-specific model & None & 0 &
Before evaluation \\
Zero-shot NAS + fine-tuning & Top-12 shortlist & Method-specific models &
Adam/Huber--50 & $\leq12$ & Before evaluation \\

H-Meta-NAS-based adaptation & Target architecture search &
Architecture-indexed meta-initializations & SGD/MSE--50 &
12.00 & Before search evaluation \\

MSA-DTI & Candidate bank & Source-trained initializations & SGD/MSE--50 & 4--7 &
Before adaptation and output \\
\bottomrule
\end{tabularx}
\vspace{0.15em}

\footnotesize
PT+FT and MSA-DTI use the same source-trained reference and target adaptation
procedure. Few-shot NAS and zero-shot NAS with fine-tuning use Adam/Huber
($\delta=0.2$), learning rate $3\times10^{-4}$, weight decay 0, and a
gradient-norm limit of 5.0. H-Meta-NAS-based adaptation uses 50 SGD/MSE
updates per adapted model in the main comparison. It uses the same 66 architecture
definitions, filters infeasible architectures before search, uses source-only
first-order meta-initializations, and selects the final architecture using
target validation data. Its settings are fixed before held-out evaluation.
All methods use the same target cases, deployment limits, and architecture
definitions.
\end{suppTableBlock}

The NAS labels in the figures and tables are abbreviated names of task-matched
controls, not exact reproductions of published NAS algorithms. The
\emph{few-shot NAS} control combines source-prior ranking with a
support-gradient entropy proxy, adapts up to 12 feasible models, and selects
by validation MSE. \emph{Zero-shot NAS} chooses one model with the
support-gradient proxy but performs no target parameter updates; here
``zero-shot'' means no adaptation updates, not no access to labeled support.
The fine-tuning variant adapts up to 12 proxy-ranked candidates before
validation selection. The H-Meta-NAS-based control instead uses separate
source meta-initializations.

\subsection{Paired held-out baseline comparisons}

Table~\ref{tab:supp-paired-baselines} summarizes the paired MSE comparisons
between MSA-DTI and the main baselines on the 80 held-out cases.

\begin{suppTableBlock}
{Paired MSE comparisons between MSA-DTI and the main baselines on the held-out cases.}
{tab:supp-paired-baselines}
\scriptsize
\renewcommand{\arraystretch}{1.08}
\setlength{\tabcolsep}{3.0pt}
\begin{tabular}{lrrr}
\toprule
Baseline
& MSE win rate (\%)
& Mean paired MSE reduction (\%)
& 95\% CI (\%) \\
\midrule
PT+FT
& 71.25 & 14.60 & $[9.41,20.14]$ \\
MeDeT-based adaptation
& 87.50 & 35.37 & $[25.39,44.26]$ \\
Random initialization
& 100.00 & 80.71 & $[78.07,83.22]$ \\
Few-shot NAS
& 95.00 & 40.40 & $[32.65,47.58]$ \\
Zero-shot NAS
& 100.00 & 76.29 & $[68.95,82.65]$ \\
Zero-shot NAS + fine-tuning
& 98.75 & 45.03 & $[36.53,52.55]$ \\
H-Meta-NAS-based adaptation
& 100.00 & 68.94 & $[64.06,73.53]$ \\
\bottomrule
\end{tabular}

\vspace{0.15em}
\footnotesize
MSE win rate compares saved test MSE with the separately evaluated baseline
using a strict inequality; ties are non-wins. For PT+FT, 13 of the 33
reference-retained cases count as strict wins because of numerical differences
of at most $7.7\times10^{-9}$ in test MSE, not additional beneficial
selections. The Figure~6 alternative labels use the released $10^{-6}$
relative-gain cutoff, which agrees with strict test-MSE ordering for all 47
selected alternatives. Positive reductions favor MSA-DTI. The 95\% intervals
are the 2.5th and 97.5th percentiles of 4000 center-cluster bootstrap
replicates, keeping each sampled target center's four cases together.
\end{suppTableBlock}

\subsection{Optimizer-matched diagnostic comparison}

This comparison examines whether the performance ordering changes when
target adaptation uses SGD/MSE and 50 parameter updates per adapted model.
All methods use the same independent diagnostic target cases and deployment
limits. The two common-shortlist controls use a shared 12-candidate set
containing A57 and the 11 highest-ranked feasible alternatives under a fixed
fused ranking (0.60 source-prior rank + 0.40 proxy rank). They differ only in
the final selection rule: lowest validation loss versus the 10\%
reference-based margin. All selections are fixed before test evaluation.

\begin{suppTableBlock}{Full optimizer-matched comparison on the diagnostic target pool.}{tab:supp-matched}
\small
\renewcommand{\arraystretch}{1.08}
\begin{tabular}{lrrrr}
\toprule
Method & MSE$\downarrow$ & Worst-10\%$\downarrow$ &
CVaR90$\downarrow$ & Adapted models \\
\midrule
MSA-DTI & \textbf{0.003809} & \textbf{0.015094} &
\textbf{0.008696} & 6.10 \\
PT+FT & 0.004497 & 0.017144 & 0.011878 & 1.00 \\
Few-shot NAS & 0.006715 & 0.025246 & 0.019653 & 12.00 \\
Zero-shot NAS + fine-tuning & 0.007402 & 0.028099 & 0.020402 & 12.00 \\
Common shortlist, lowest validation loss &
0.006707 & 0.025240 & 0.019653 & 12.00 \\
Common shortlist, reference-based selection &
0.006819 & 0.025562 & 0.020547 & 12.00 \\
\bottomrule
\end{tabular}
\end{suppTableBlock}

MSA-DTI gives the lowest MSE, Worst-10\% error, and CVaR90 in this
comparison. Source initializations, selection rules, and numbers of adapted
models remain method-specific, so the common per-model budget does not imply
equal total target-side work.

\subsection{Selection-outcome details}

On the 80 independent diagnostic cases, the relative-margin rule selects
alternatives in 61 cases, including 49 beneficial and 12 harmful selections.
Direct lowest-validation-loss selection chooses alternatives in 69 cases,
including 55 beneficial and 14 harmful selections. The corresponding
all-case harmful-selection rates are 15.00\% and 17.50\%, respectively.


\paperfig{0.88\linewidth}{fig9.pdf}
{Selection outcomes on the 80-case diagnostic pool.}
{fig:supp-selection-diagnostic}

The relative-margin rule retains the reference more often. The individual
harmful-selection-rate intervals overlap, but their overlap does not test the
paired difference. The observed rate difference is reported descriptively.

\section{Timing, Hardware Profiling, and Deployment Analysis}
\label{supp:hardware}

\subsection{Target-side timing protocol}

All methods are implemented in Python and PyTorch and run on an NVIDIA RTX
3060 Laptop GPU. Target-side time includes model creation, initialization
loading, target-side search when applicable, target adaptation, validation
evaluation, and final selection. It excludes offline preparation and test
evaluation. Each reported time is the mean over five runs, with the
corresponding standard deviation reported in the main manuscript.

\subsection{Architecture-profiling protocol}

The five executable architectures are profiled with batch size one on an
Intel Core i7-12700H CPU and an NVIDIA RTX 3060 Laptop GPU for
$H\in\{1,4\}$. Each setting uses 100 warm-up inferences followed by five
repetitions of 1000 timed inferences. CUDA execution is synchronized before
and after each timed region. Median latency is the primary latency measure,
and the 95th percentile summarizes upper-tail latency. Serialized
parameter-file size is used as the measured model size. The increment in
peak allocated CUDA tensor memory is recorded for each architecture--horizon
pair.

\begin{suppTableBlock}{Per-architecture batch-one hardware profile.}{tab:supp-hardware-profile}
\scriptsize
\renewcommand{\arraystretch}{1.05}
\setlength{\tabcolsep}{2.0pt}
\resizebox{\linewidth}{!}{%
\begin{tabular}{lrrrrrrr}
\toprule
Configuration & $H$ & Operations & Parameters & CPU (ms) & GPU (ms) &
State (KiB) & GPU $\Delta$peak (KiB) \\
\midrule
A1  & 1 & 155{,}712   & 77{,}921 & 0.097 & 0.148 & 306.5 & 1024 \\
A1  & 4 & 155{,}904   & 78{,}020 & 0.094 & 0.133 & 306.9 & 1024 \\
A6  & 1 & 157{,}760   & 78{,}977 & 0.093 & 0.168 & 311.0 & 1024 \\
A6  & 4 & 157{,}952   & 79{,}076 & 0.099 & 0.151 & 311.4 & 1024 \\
A13 & 1 & 159{,}808   & 80{,}033 & 0.128 & 0.237 & 315.5 & 1024 \\
A13 & 4 & 160{,}000   & 80{,}132 & 0.128 & 0.218 & 315.9 & 1024 \\
A55 & 1 & 377{,}888   & 2{,}081  & 1.793 & 3.999 & 10.3  & 54 \\
A55 & 4 & 377{,}984   & 2{,}132  & 1.793 & 4.004 & 10.4  & 1030 \\
A56/A57 & 1 & 1{,}050{,}688 & 5{,}697 & 1.980 & 4.019 & 24.4 & 60 \\
A56/A57 & 4 & 1{,}050{,}880 & 5{,}796 & 1.975 & 4.170 & 24.8 & 1037 \\
\bottomrule
\end{tabular}}
\vspace{0.15em}

\footnotesize\raggedright
CPU and GPU values are median inference latency. A56 and A57 are the same
executable one-layer GRU-32 architecture. GPU $\Delta$peak is the
nonnegative difference between the PyTorch CUDA maximum allocated tensor
memory during the timed repetitions and the allocated-memory baseline
recorded after warm-up and a cache clear, with peak statistics reset before
timing. It excludes resident model/input allocations and is neither total
device usage nor a configured deployment memory limit.
\end{suppTableBlock}

\subsection{Profiling results}

For each device and $H\in\{1,4\}$, the architecture-level correlation
uses the five profiles displayed in Table~\ref{tab:supp-hardware-profile},
with one representative entry for the executable A56/A57 architecture.
Across these four device--horizon groups ($n=5$ distinct architectures per
group), the Spearman correlation between estimated operation count and
median inference latency ranges from 0.900 to 1.000. The correlation between
parameter count and serialized parameter-file size is 1.000 in every measured setting, supporting
parameter count as a model-size indicator. Peak GPU allocation is not
monotonic in parameter count across the two horizons, so parameter count is
not a direct proxy for runtime GPU memory. The MLP architectures use fewer
estimated operations but more parameters than the GRU architectures. The two
deployment indicators therefore capture different model properties, while
their activity depends on the candidate set and limit tier. The profiling
results characterize only the measured CPU/GPU host.

\subsection{Single-limit feasible-space analysis}

To identify the active constraint in the main candidate bank, we apply the
operation-count and parameter-count limits separately to the same seven
candidates using their fixed complexity profiles.

For both $H=1$ and $H=4$, all seven candidates satisfy the operation-count
limit under the tight, medium, and loose settings. Under the tight
parameter-count limit, the three MLP candidates are excluded and the
feasible bank is reduced from seven candidates to four, whereas the medium
and loose parameter-count limits retain all seven candidates. Consequently,
the joint feasible set contains four candidates under the tight setting and
seven under the medium and loose settings. The fixed reference candidate
satisfies both limits in all six horizon--tier combinations.

\subsection{Expanded-bank dual-limit stress analysis}

The preceding analysis shows that the operation-count limit is inactive for
the seven-candidate main bank under the main limit tiers.
To examine whether operation count and parameter count can independently
restrict candidate feasibility, we additionally construct an eight-candidate
stress bank by adding configuration 24 from the predefined architecture
space.

The stress thresholds are derived solely from deterministic model-complexity
profiles. Across $H\in\{1,4\}$, the largest operation count in the
original seven-candidate bank is $1{,}050{,}880$, whereas the smallest
operation count of configuration 24 is $1{,}198{,}112$. Their midpoint gives
the stress operation-count limit $B^{F}=1{,}124{,}496$. Similarly, the
largest parameter count among the stress-bank candidates other than A13 is
$79{,}076$, whereas the smallest parameter count of A13 is $80{,}033$.
Their midpoint gives the stress parameter-count limit
$B^{P}=79{,}554.5$.

For both $H=1$ and $H=4$, configuration 24 violates only the
operation-count limit, A13 violates only the parameter-count limit, and the
fixed reference A57 satisfies both limits. Each single-limit feasible set
therefore retains seven of the eight candidates, whereas their joint
feasible set retains six. Thus, the joint feasible set is a strict subset
of each single-limit feasible set.

The target-side stress evaluation uses the same 80 held-out cases, and all
models selected under the joint limits satisfy both limits. The expanded bank
is used only for this limit-complementarity diagnostic, whereas the main
evaluation uses the original seven-candidate bank.


\subsection{Selected-architecture analysis}

\paperfig{0.88\linewidth}{fig10.pdf}
{Selected-configuration complexity and performance. Marker area shows the
number of selected cases, and $h$ denotes harmful selections.}
{fig:supp-architecture-tradeoff}

The reference configuration is retained in 33 cases. Four alternative
configurations cover the remaining 47 cases, and their mean paired MSE
reductions relative to the independently evaluated PT+FT baseline range
from 17.43\% to 30.50\%. All four alternatives have positive mean
reductions. The three harmful selections occur on two alternative GRU configurations.

\section{Configuration Analysis}
\label{supp:configuration}

\paperfig{0.92\linewidth}{fig11.pdf}
{Configuration analysis of (a) the retained configuration count,
(b) the target update budget, and (c) the relative margin.}
{fig:supp-configuration}

Diagnostic MSE decreases from 0.004917 with one retained configuration to
0.004063 with four retained configurations and then remains unchanged. The
main evaluation uses the six-configuration list fixed before evaluation.
The observed saturation after four configurations suggests diminishing
returns from adding retained configurations in this diagnostic; it does not
identify an optimal bank size.

Most of the observed reduction in diagnostic check MSE occurs within the
first ten target updates. The main experiment nevertheless uses 50 updates
for every feasible candidate to preserve a common fixed target budget.

Using the same candidate bank, source initializations, feasibility rule, and
reference-based selection rule at both endpoints, selected check
MSE is 0.003973 with zero target updates and 0.003784 after 50 updates. The relative reduction computed from the two aggregate mean
check-MSE values is 4.75\%. This is a ratio of means, not the mean of
paired case-wise relative reductions. The same candidate is selected at
both endpoints in 51.25\% of cases. An earlier interval for a different
paired estimand cannot be assigned to this aggregate ratio without its
original calculation record; accordingly, no interval-based benefit claim
is made for this diagnostic.

The relative-margin analysis shows that 10\% is the smallest value satisfying
all preset calibration criteria.

\section{Source and Preselection Sensitivity and Real Workloads}
\label{supp:external}

\subsection{Source-center scale and source-training seeds}

\paperfig{0.92\linewidth}{fig12.pdf}
{Synthetic sensitivity to (a) source-center scale and (b) source-training
seeds.}
{fig:supp-source-sensitivity}

The mean paired MSE reductions relative to the matched reference at
source-center scales 10, 20, 30, 40, and 50 are 6.46\%, 4.92\%, 3.78\%,
3.41\%, and 4.14\%, respectively. The mean gain remains positive at every
tested source scale.

Across source-training seeds 2904, 2905, and 2906, the mean paired MSE
reductions are 10.52\%, 13.21\%, and 11.25\%. All three confidence
intervals remain above zero, although some individual cases have zero or
negative gains.
The retained architecture list is held fixed in this seed study. The result
therefore measures sensitivity to source-model training and initialization.

\subsection{Independent architecture-preselection stability}

We separately repeat the complete 66-configuration diagnostic preselection
on two new 20-center pools. Together with the original 20-center preselection
pool, this gives three independently sampled preselection sets. The A57
reference, architecture space, initialization recipe, SGD/MSE--50 target
recipe, and the preselection and retention rules defined in Supplementary
Section~S2.1 are kept unchanged. None of these pools overlaps the 20-center
robustness pool used below, and no test data are used during preselection.

\begin{suppTableBlock}{Independent architecture-preselection stability and evaluation on centers 1160--1179.}{tab:supp-screening-stability}
\scriptsize
\renewcommand{\arraystretch}{1.08}
\setlength{\tabcolsep}{2.2pt}
\resizebox{\linewidth}{!}{%
\begin{tabular}{llllrrrr}
\toprule
Preselection pool & Centers & Retained configurations & Size & Jaccard &
Gain (\%) & 95\% CI (\%) & Harmful (\%) \\
\midrule
Main S0 & 900--919 & A1, A6, A13, A55, A56, A57 & 6 & 1.000 & 0.295 & $[-0.646,1.564]$ & 2.50 \\
Independent S1 & 1120--1139 & reference fallback only & 0 & 0.000 & 0.000 & $[0.000,0.000]$ & 0.00 \\
Independent S2 & 1140--1159 & A55, A56, A57, A59 & 4 & 0.429 & 0.125 & $[-0.894,1.427]$ & 3.75 \\
\bottomrule
\end{tabular}}
\vspace{0.15em}

\footnotesize
Jaccard similarity and Size are computed on the indexed-configuration lists retained by preselection, including A57 whenever it occurs in a list.
The PT-A57 reference candidate is added separately to every downstream
candidate bank; a retained A57 additionally supplies the paired C1-A57
candidate. Gain is the mean paired check-MSE reduction relative to PT-A57 on
80 robustness cases; intervals use a 4000-repetition center-cluster bootstrap.
The diagnostic evaluates preselection stability while keeping the downstream
procedure unchanged.
\end{suppTableBlock}

Across S0--S2, retained-list sizes are 6, 0, and 4 (median 4, range 0--6).
A55, A56, and A57 are retained in two of three pools; A1, A6, A13, and A59
are each retained once. On the robustness pool, the S0 list selects PT-A57
in 72 cases, A55 in three, A56 in three, and the paired-control A57
initialization in two. S1 falls back to PT-A57 in all 80 cases. S2 selects
PT-A57 in 71 cases, A55 in three, A56 in three, the paired-control A57
initialization in two, and A59 in one.

The two nonempty retained lists, S0 and S2, both have 95\% confidence
intervals for paired check-MSE gains relative to the adapted PT-A57 reference
that include zero on the robustness pool, indicating sample sensitivity in
candidate discovery. The main evaluation uses the S0 bank obtained from the
main-study preselection pool, while the A57 reference provides a fallback when a
preselection pool retains no alternative.

Reapplying the recorded preselection counts with retention thresholds of one,
two, and three wins retains 10, 6, and 2 configurations, respectively. The threshold-two list used in the main study is S0, whereas threshold three retains only A55 and
A57. On the existing step-50 diagnostic trajectory, threshold three increases
selected check MSE by 3.79\% relative to threshold two, with a 95\% confidence
interval of $[1.58\%,6.30\%]$. The threshold-one list additionally contains A4, A5, A12, and A59. The main study did not generate source-trained records for these four
additional configurations; downstream threshold sensitivity is therefore
reported only for thresholds two and three.

\subsection{End-to-end propagation of preselection outcomes}

To characterize the downstream effect of candidate-discovery variability,
the recorded S0--S2 preselection outcomes are treated as fixed inputs and carried
through the complete downstream procedure. Each outcome uses the same source
centers and source-training recipe, followed by bank-specific application of the same margin-calibration rule on the same disjoint 20-center
calibration pool. The resulting pipelines are evaluated on a new 20-center
held-out pool (centers 1180--1199; 80 cases). The architecture space, A57
reference, deployment limits, target adaptation protocol, margin grid, and
calibration criteria remain unchanged. Held-out test data are not used in
source training, calibration, deployment filtering, adaptation, or selection.
For S2, A59 is source-trained using the same source-training recipe; S1
contains only the PT-A57 reference candidate.

\begin{suppTableBlock}
{End-to-end propagation of the three preselection outcomes on centers
1180--1199.}
{tab:supp-screening-e2e}
\footnotesize
\renewcommand{\arraystretch}{1.12}
\setlength{\tabcolsep}{2.0pt}
\textit{(a) Candidate banks and calibrated margins}\par\smallskip
\begin{tabularx}{\linewidth}{Y Y Z{0.11\linewidth} Z{0.11\linewidth}}
\toprule
Preselection outcome & Retained configurations & Bank size & $\tau$ \\
\midrule
S0 & A1, A6, A13, A55, A56, A57 & 7 & 0.075 \\
S1 & Reference fallback only & 1 & -- \\
S2 & A55, A56, A57, A59 & 5 & 0.125 \\
\bottomrule
\end{tabularx}
\medskip
\textit{(b) Held-out performance and selection outcomes}\par\smallskip
\begin{tabularx}{\linewidth}{Y
Z{0.085\linewidth} Z{0.215\linewidth} Z{0.10\linewidth}
Z{0.09\linewidth} Z{0.16\linewidth} Z{0.14\linewidth}}
\toprule
Preselection outcome & Gain (\%) & 95\% CI (\%) & MSE win (\%) &
Alt. selected & Harmful (\%) & Limit satisfaction (\%) \\
\midrule
S0 & 10.034 & $[4.602,16.019]$ & 52.50 & 48 & 7.50 / 12.50 & 100.00 \\
S1 & 0.000 & $[0.000,0.000]$ & 0.00 & 0 & 0.00 / -- & 100.00 \\
S2 & 8.784 & $[3.833,14.238]$ & 41.25 & 39 & 7.50 / 15.38 & 100.00 \\
\bottomrule
\end{tabularx}
\vspace{0.15em}

\footnotesize
Gain is the mean paired test-MSE reduction relative to PT+FT. The 95\%
intervals use 4000 center-cluster bootstrap repetitions with the four cases
of each target center kept together. Harmful rates are reported as
all-case / conditional-on-alternative-selection percentages. For S1, no
alternative is selected, so the conditional rate is undefined and shown as
``--''. The adapted PT-A57 reference is used in all S1 cases.
\end{suppTableBlock}

The two nonempty banks yield positive mean paired MSE reductions with
confidence intervals above zero despite their different retained
configurations. The S1 outcome contains no alternative and therefore follows
the adapted-reference fallback without changing performance relative to
PT+FT. These end-to-end results characterize the downstream consequence of
the observed preselection variability; the retained-bank composition itself
remains sample-sensitive.

\subsection{Causally corrected Alibaba workload evaluation}

We use Alibaba Cluster Trace v2018 in a reconstructed causal-correction
evaluation. The 20 source machines are deterministically reconstructed, and
the 20 calibration and 40 held-out machine identities match the retained
historical cohorts; their roles remain disjoint. The corrected source models
are retrained on an independent source pool. A source-derived 10-second event-
time grid supplies the fixed prediction opportunities, with event time assumed
to be the observation-availability time. The input length is $L=96$ and the
horizon is $H\in\{1,4\}$. The response is the next $H$ genuinely observed
CPU values scaled from 0--100 to fractions in $[0,1]$, so MSE has units of
CPU-fraction squared. The 25-channel input uses historical values,
forward-only missing-value handling, trailing-six CPU statistics, and causal
observation masks. Source MKPI values were unavailable in the registered
source interval; its value and mask channels are identically zero for all
methods. This correction therefore does not establish performance with MKPI.

For each target case, adaptation uses $K\in\{10,20\}$ valid labeled support
windows, followed by 80 distinct validation windows; the calibration cases
add 60 check windows. Held-out testing evaluates every label-valid origin
within 320 fixed candidate opportunities (290--319 windows per case).
Missing labels are not interpolated, and every method uses the same eligible
evaluation origins. Deployment-limit tiers remain controlled settings rather
than machine-measured resource limits. The main and paired selectors are
based on separately source-trained initializations, with the paired control
selecting between the two A57 candidates. Their source-training procedures
are not identical, so this is not a seed-only control.

The full five-criterion rule of main Eq.~(12) is applied to all six fixed
margins on the 80 calibration cases, with a 20-machine, 4,000-repetition
paired cluster bootstrap (seed 2904). The reference and paired gains are
case-wise relative reductions in check-set MSE. Table~\ref{tab:supp-alibaba-calibration}
reports the corresponding mean gains and 95\% lower confidence bounds.

\begin{suppTableBlock}{Full five-criterion calibration for the causally corrected Alibaba cases (20 machines, 80 cases).}{tab:supp-alibaba-calibration}
\footnotesize
\setlength{\tabcolsep}{4.0pt}
\renewcommand{\arraystretch}{1.1}
\begin{tabular}{lrrrrrl}
\toprule
$\tau$ & Harm (\%) & Ref. gain (\%) & Ref. LCB (\%) & Pair gain (\%) & Pair LCB (\%) & Eligible\\
\midrule
0.050 & 15.00 & -6.76 & -15.06 & -2.05 & -8.51 & No\\
0.075 & 15.00 & -6.96 & -15.29 & -2.25 & -8.63 & No\\
0.100 & 8.75 & -5.68 & -13.70 & -3.33 & -10.18 & No\\
0.125 & 8.75 & -5.68 & -13.70 & -3.33 & -10.18 & No\\
0.150 & 6.25 & -4.11 & -11.04 & -2.28 & -8.41 & No\\
0.200 & 2.50 & -1.43 & -6.98 & -1.43 & -6.98 & No\\
\bottomrule
\end{tabular}
\vspace{0.3em}
\par\footnotesize\raggedright
Harm is the fraction of all calibration cases with a selected alternative
whose check loss exceeds the reference. Eligibility requires harm $\le5\%$,
both mean gains $\ge3\%$, and both respective 95\% lower bounds $>0$.
All six margins fail the complete rule.
\end{suppTableBlock}

Consequently, the calibrated MSA-DTI policy retains the adapted PT-A57
reference for all 160 held-out cases. It has the same mean MSE (0.00338831)
as PT+FT and an identical-case paired improvement of 0 by construction;
this is not evidence of a population-level harmful-selection guarantee.
The separately pre-registered fixed-$\tau=0.10$ diagnostic selects nine
alternatives, including three harmful replacements (3/160, or 1.875\% of
all held-out cases). The mean of case-wise relative MSE gains against PT+FT
is $-0.2592\%$ (machine-cluster 95\% CI $[-1.4711\%,0.6205\%]$).
Although its mean MSE (0.00337210) is slightly below the PT+FT mean,
that ratio of mean MSEs is not the reported mean paired relative gain.
The diagnostic does not establish positive external-domain transfer.
\paperfig{0.95\linewidth}{fig_alibaba_corrected_heldout.pdf}
{Causally corrected Alibaba held-out comparisons against PT+FT: (a) case-level
relative-gain distributions and (b) machine-cluster 95\% confidence
intervals for mean case-wise relative gains. Fixed-margin results are
separate diagnostics; the calibrated policy retains the reference.}
{fig:supp-alibaba-causal-heldout}

The corrected evaluation uses a previously observed machine population;
it is a post-observation correction, not an independent blind study. Arrival times were not recorded and event timestamps are assumed to
be availability times. The earlier noncausal Alibaba result is preserved as
historical evidence but is not used to support the present external-domain
claims. A separate corrected reproducibility archive retains protocol locks,
model checkpoints, case-level records, and plots; it has not yet been
publicly deposited. The third-party raw trace and large corrected input/label
arrays are not redistributed.


\subsection{Bitbrains in-domain VM workload evaluation}
\label{supp:bitbrains}

The publicly released GWA-T-12 Bitbrains fastStorage trace\footnote{Dataset: \url{https://atlarge-research.com/gwa-t-12/}.
The trace was provided by Bitbrains IT Services Inc.} contains native
VM-level CPU, memory, storage, and network observations. A registered quality
filter was applied to the full set of 1,250 VMs, after which 100 eligible VMs
were selected deterministically by their stable identifiers. Disjoint sets of
20 source, 20 development, 20 calibration, and 40 held-out VMs were fixed
before model evaluation. This is an in-domain model-instantiation experiment:
source initializations are trained anew on Bitbrains source VMs rather than
transferred as frozen synthetic or Alibaba weights. The prediction target is
recorded CPU usage in percent divided by 100; values above 100\% are retained
as recorded, not clipped. The input uses source-only normalization, causal
historical features, and role-disjoint chronological windows; $H=1,4$ means
next one or four native records. Typical elapsed horizons are about 300 and
1200 seconds, although irregular record gaps cause variation.

Within each held-out case, MSA-DTI and each comparator use the same observed
Test targets. Support adaptation uses $K=10$ or $20$ windows per model with
50 SGD/MSE updates. The pretrained seven-token bank, per-model feasibility
checks, and the five relative-margin criteria remain as in the synthetic
method; calibration is performed on distinct VMs. For the development-fixed
control, one feasible candidate is selected in advance for each budget tier
and horizon using average development check MSE. It alone is adapted at the
target; it does not inspect target validation for architecture selection.
The full MSA-DTI selector instead adapts the 4 or 7 feasible candidates
before comparing their post-adaptation validation scores. Thus, the method
and fixed control share the candidate bank and target protocol but not the
total number of candidate adaptations.

\begin{suppTableBlock}{Full five-criterion relative-margin calibration on Bitbrains development-disjoint calibration VMs (20 VMs, 80 cases).}{tab:supp-bitbrains-calibration}
\footnotesize
\renewcommand{\arraystretch}{1.09}
\setlength{\tabcolsep}{3pt}
\begin{tabularx}{\linewidth}{
Z{0.10\linewidth}
Z{0.12\linewidth}
Y
Y
Z{0.13\linewidth}}
\toprule
$\tau$ & Harm/all cases & Gain vs. reference (95\% CI) & Gain vs. paired selector (95\% CI) & Eligible \\
\midrule
0.050 & 0.00\% & 38.56\% [25.64\%,51.83\%] & 19.36\% [10.30\%,29.46\%] & Yes \\
0.075 & 0.00\% & 38.39\% [25.37\%,51.65\%] & 19.39\% [10.26\%,29.47\%] & Yes \\
0.100 & 0.00\% & 38.01\% [24.87\%,51.40\%] & 19.01\% [9.96\%,29.03\%] & Yes \\
0.125 & 0.00\% & 38.01\% [24.87\%,51.40\%] & 18.98\% [9.95\%,29.01\%] & Yes \\
0.150 & 0.00\% & 37.81\% [24.59\%,51.31\%] & 18.97\% [9.95\%,28.98\%] & Yes \\
0.200 & 0.00\% & 37.81\% [24.59\%,51.31\%] & 18.97\% [9.95\%,28.98\%] & Yes \\
\bottomrule
\end{tabularx}
\vspace{0.15em}
\footnotesize All rows pass all five criteria, including the preset 5\% all-case harmful rate,
both minimum 3\% gains, and both strictly positive 95\% lower bounds.
The least restrictive qualifying margin, $\tau=0.05$, is frozen before held-out evaluation.
Intervals use 4,000 VM-cluster bootstrap repetitions. These empirical
calibration conditions are not a held-out or population guarantee.
\end{suppTableBlock}

For the 160 held-out cases, MSA-DTI achieves a 32.7045\% mean of case-wise
relative test-MSE improvements over PT+FT (95\% VM-cluster CI
[7.6721\%,51.2506\%]); 123 cases improve, 11 worsen, and 26 are equal.
The mean case MSE is 0.00742182 for MSA-DTI versus 0.00807491 for PT+FT.
Their ratio of aggregate mean losses corresponds to an 8.0879\% reduction;
this secondary estimand is not interchangeable with the mean of case ratios.
The comparison against the development-fixed control yields a descriptive
38.3851\% mean case-relative gain (95\% VM-cluster CI
[25.5620\%,50.8000\%]). The held-out harmful-selection frequency is
11/160 (6.8750\%), exceeding the 5\% empirical calibration criterion.
The primary evidence supports improvement in this VM cohort but not universal
risk control or benefit on all real traces, as the corrected Alibaba result
demonstrates.

\paperfig{0.94\linewidth}{fig_bitbrains_gain.pdf}
{Bitbrains fastStorage case-relative held-out MSE gains for calibrated
MSA-DTI versus PT+FT. Intervals use 4,000 paired VM-cluster bootstrap
draws; subgroup comparisons are descriptive without multiplicity correction.
All registered cases, including large adverse relative values, are retained.}
{fig:supp-bitbrains-gain}

The mean target-side ready-tensor processing time is 19.4735 seconds per
MSA-DTI case versus 4.7403 seconds for PT+FT, with mean numbers of
adapted models 6.10 versus 1.00. Shared execution components limit strict
interpretation as independent end-to-end timing runs. Hidden VM host/tenant
structure, a single source-training realization, and assumed event-time
data availability also limit broader generalization. The primary fastStorage
cohort was selected before model outcomes; the optional Rnd trace was not
run in this study and is not an unreported negative or positive result.

\subsection{Strong development-fixed candidate comparison on synthetic centers}
\label{supp:r5-fixed}

A targeted supplementary contrast evaluates whether target-specific
post-adaptation selection improves on committing in advance to a strong
candidate from the same seven-token source-trained bank. For each budget
and horizon, an independent development-center pool selects the lowest
mean-check-MSE feasible token after averaging the two support sizes per
center; the chosen token is frozen before the 80 evaluation cases. During
evaluation, the fixed comparator adapts only that token for 50 updates,
whereas MSA-DTI adapts all feasible tokens under the same per-model update
rule and uses the fixed synthetic margin $\tau=0.10$ for validation
selection. The fixed comparator therefore has lower total target effort.

The MSA-DTI mean case-relative test-MSE improvement over this
\emph{Dev-Fixed-Best} comparator is 13.4803\% (95\% center-cluster CI
[6.3856\%,21.0963\%]); 43 cases improve, 12 worsen, and 25 are equal.
MSA-DTI requires 5.95 adapted models and 297.5 updates per case on average,
versus 1 model and 50 updates for Dev-Fixed-Best. Mean ready-tensor
target-side time is 15.5813 versus 1.8262 seconds, respectively.
\emph{Evidence-status limitation:} the new development and evaluation pools
do not overlap any reconstructable historical center/seed domains, but the
original RNG identities of two historical R2 screening pools are not
recoverable. This comparison is retained as post-audit supplementary
support, not as strictly confirmed independent evidence closing every
historical development-pool separation requirement.

\paperfig{0.94\linewidth}{fig_r5_devfixed_gain.pdf}
{Additional synthetic same-bank strong-fixed control: case-relative test-MSE
gains of MSA-DTI over Dev-Fixed-Best. Paired intervals resample 20 centers;
historical development-pool RNG provenance is not fully closed.}
{fig:supp-r5-fixed-gain}

\section{Synthetic Benchmark Generation and Center Heterogeneity}
\label{supp:synthetic-benchmark}

The controlled benchmark represents heterogeneous computing centers in a
common metric coordinate space.
Each center may observe only a subset of the metric positions.
The model input uses 12 value channels, 12 corresponding observation-mask
channels, and one time-gap channel, giving a fixed input dimension of 25.
Unavailable metric values are zero-filled and identified by the
corresponding mask channels.

Center types A--C are evaluation strata rather than model classes.
They differ in workload dynamics, sampling, and monitoring conditions.
Deployment-limit tiers are separate platform-side settings and are assigned
using the type-dependent probabilities reported in
Table~\ref{tab:supp-center-types}.
The table also gives the complete type-specific parameter ranges.
The default A:B:C ratio is 3:4:3, giving 6 Type-A, 8 Type-B, and 6 Type-C
centers in each 20-center source or held-out target pool.

\begin{suppTableBlock}
{Type-specific parameter ranges of the controlled synthetic benchmark.}
{tab:supp-center-types}
\scriptsize
\renewcommand{\arraystretch}{1.06}
\setlength{\tabcolsep}{3.2pt}
\begin{tabularx}{0.98\linewidth}{
    Y
    Z{0.21\linewidth}
    Z{0.21\linewidth}
    Z{0.21\linewidth}}
\toprule
Parameter & Type A & Type B & Type C \\
\midrule

Sampling interval $d_t$
& $\{1\}$
& $\{1,2\}$
& $\{2,3,4\}$ \\

Point-missing rate
& $[0.02,0.08]$
& $[0.10,0.25]$
& $[0.25,0.45]$ \\

Missing-block length
& $[5,15]$
& $[20,60]$
& $[60,120]$ \\

Observation-noise std.
& $[0.01,0.03]$
& $[0.03,0.06]$
& $[0.06,0.10]$ \\

Permanently unavailable metrics
& $0$--$2$
& $2$--$4$
& $4$--$6$ \\

Base workload logit
& $[-0.70,0.10]$
& $[-0.45,0.35]$
& $[-0.20,0.55]$ \\

Seasonal amplitude
& $[0.20,0.45]$
& $[0.30,0.60]$
& $[0.40,0.75]$ \\

AR coefficient
& $[0.55,0.80]$
& $[0.60,0.88]$
& $[0.65,0.92]$ \\

Process-noise std.
& $[0.04,0.08]$
& $[0.06,0.11]$
& $[0.08,0.14]$ \\

Burst rate per 1000 base steps
& $[1.0,2.5]$
& $[1.5,3.5]$
& $[2.0,4.5]$ \\

Burst amplitude
& $[0.25,0.55]$
& $[0.35,0.75]$
& $[0.45,0.95]$ \\

Drift probability
& $0.20$
& $0.40$
& $0.60$ \\

Drift strength
& $[0.10,0.25]$
& $[0.15,0.35]$
& $[0.20,0.45]$ \\

Limit-tier probabilities
(tight / medium / loose)
& $0.25/0.45/0.30$
& $0.35/0.45/0.20$
& $0.45/0.40/0.15$ \\

\bottomrule
\end{tabularx}
\end{suppTableBlock}

The three types share the same missingness mechanisms.
For every center, the generator may create 2--6 single-metric outage
blocks and 1--3 grouped outage blocks containing 2--4 metrics.
A short global monitoring outage occurs with probability 0.20.
The shared periodic base scales are 24, 96, and 672 base steps.

Together, these settings vary five aspects of the benchmark: workload
dynamics, sampling, monitored-variable availability, observation conditions,
and deployment-limit settings.

For each center, the generator first creates a latent workload process.
The workload contains periodic variation, autoregressive variation, and
random workload bursts.
When drift occurs, it is generated as either an abrupt change or a gradual
transition.

The latent workload is then mapped to multivariate monitoring signals.
Individual metrics may have different delays, gains, offsets, nonlinear
terms, and observation noise.
Center-specific sampling is applied before monitoring-schema and missingness
effects are introduced.

The benchmark generator and all data-role assignments are fixed before
architecture preselection, margin calibration, diagnostics, and held-out
evaluation.

\end{document}
